\subsection{FINER STRUCTURES}

For the rest of this section we shall suppose that we are given a species of structures \(\Sigma\) and a notion of \(\sigma\) -morphism relative to this species of structures; all the notions which will be introduced will depend not only on \(\Sigma\) but also on the notion of \(\sigma\) -morphism envisaged. Usually we shall say ``morphism'' in place of `` \(\sigma\) -morphism''.

Let E be a set and let \(S_{1}\) , \(S_{2}\) be two structures of species \(\Sigma\) on E. The structure \(S_{1}\) is said to be finer than \(S_{2}\) (and \(S_{2}\) coarser than \(S_{1}\) ) if the identity mapping of E, endowed with \(S_{1}\) , onto E, endowed with \(S_{2}\) , is a morphism.

When necessary to avoid ambiguity, we shall say that \(G_{1}\) is finer than \(G_{2}\) relative to the notion of \(\sigma\) -morphism under consideration; and similarly for all the other notions to be defined in this section.

Suppose that \(\mathcal{S}_1\) is finer than \(\mathcal{S}_2\) . If \(\mathbf{E}'\) is a set endowed with a structure \(\mathcal{S}'\) of species \(\Sigma\) , and if \(f\) is a morphism of \(\mathbf{E}\) , endowed with \(\mathcal{S}_2\) , into \(\mathbf{E}'\) , endowed with \(\mathcal{S}'\) , then \(f\) is also a morphism of \(\mathbf{E}\) , endowed with \(\mathcal{S}_1\) , into \(\mathbf{E}'\) , endowed with \(\mathcal{S}'\) ; this follows from the preceding definition and from \((\mathrm{MO}_{\mathrm{II}})\) . Likewise, if \(g\) is a morphism of \(\mathbf{E}'\) , endowed with \(\mathcal{S}'\) , into \(\mathbf{E}\) , endowed with \(\mathcal{S}_1\) , then \(g\) is also a morphism of \(\mathbf{E}'\) , endowed with \(\mathcal{S}'\) , into \(\mathbf{E}\) , endowed with \(\mathcal{S}_2\) .

Thus the finer the structure (of species \(\Sigma\) ) on E, the more morphisms there are with E as source, and the fewer morphisms with E as target.

The relation `` \(\mathcal{S}_{1}\) is coarser than \(\mathcal{S}_{2}\) '' is an order relation between \(\mathcal{S}_{1}\) and \(\mathcal{S}_{2}\) on the set of all structures of the species \(\Sigma\) on E; for it is reflexive by (MO \(_{\text{III}}\) ), transitive by (MO \(_{\text{II}}\) ), and if a structure of species \(\Sigma\) is both finer and coarser than another, then the two structures are identical by virtue of (MO \(_{\text{III}}\) ). In conformity with the general definitions (Chapter III, § 1, no. 14), two structures of species \(\Sigma\) on E are said to be comparable if one is finer than the other; a structure is said to be strictly finer (resp. strictly coarser) than another if it is finer (resp. coarser) than the other and is distinct from it.

Examples

\begin{enumerate}
\def\labelenumi{(\arabic{enumi})}
\item
  An order structure with graph \(s\) on a set \(A\) is finer than an order structure with graph \(s'\) if and only if \(s \subset s'\) . In other words, the relation \(x \leqslant y\) with respect to \(s\) implies \(x \leqslant y\) with respect to \(s'\) ; this is in accordance with the definition given in Chapter III, § 1, no. 4, Example 3.
\item
  Consider two algebraic structures F, \(F'\) of the same species \(\Sigma\) on a set A, where F and \(F'\) are the graphs of the (everywhere defined) laws of composition of these two structures. From the definition of morphisms in this case (no. 1, Example 2), F is finer than \(F'\) if and only if \(F \subset F'\) . But since F and \(F'\) are both functional graphs with the same domain \(A \times A\) , we must have \(F = F'\) . In other words, two comparable structures of species \(\Sigma\) are necessarily identical.
\item
  Let V, \(V'\) be two topologies on the same set A. To say that V is finer than \(V'\) means, by virtue of the definition of morphisms in this case (no. 1, Example 3), that \(V' \subset V\) ; in other words, every subset of A which is an open set in the topology \(V'\) is also an open set in the topology V (and thus, the finer the topology, the more open sets there are).
\end{enumerate}

Remark. We have just met an example (Example 2) in which two comparable structures of the same species \(\Sigma\) are necessarily identical. There are many other such examples: linear order structures, * compact topologies, Fréchet space structures (the morphisms being continuous linear mappings), topologies defined by an absolute value (or a valuation) on a division ring, etc. *

For such a species of structures \(\Sigma\) , a morphism \(f\) of E into E' which is bijective is an isomorphism; for if we transport to E' the structure \(\mathcal{S}\) on E by means of \(f\) , we obtain a structure of species \(\Sigma\) which is finer than the structure \(\mathcal{S}'\) on E' and therefore coincides with \(\mathcal{S}'\) .

